```latex
\documentclass[11pt]{amsart}

\usepackage{amsmath,amssymb,amsthm,mathtools}
\usepackage[margin=1.1in]{geometry}

\newtheorem{theorem}{Theorem}[section]
\newtheorem{lemma}[theorem]{Lemma}
\newtheorem{proposition}[theorem]{Proposition}
\newtheorem{corollary}[theorem]{Corollary}
\newtheorem{definition}[theorem]{Definition}
\newtheorem{remark}[theorem]{Remark}
\newtheorem{assumption}[theorem]{Assumption}
\newtheorem{example}[theorem]{Example}

\newcommand{\R}{\mathbb{R}}
\newcommand{\E}{\mathbb{E}}
\newcommand{\W}{\mathbb{W}}
\newcommand{\1}{\mathbf{1}}
\newcommand{\F}{\mathcal{F}}
\newcommand{\N}{\mathcal{N}}
\newcommand{\dd}{\mathrm{d}}
\newcommand{\divop}{\operatorname{div}}
\newcommand{\KL}{\mathrm{KL}}
\newcommand{\Lto}{\longrightarrow}

\begin{document}

\title{Minimax Estimation of a Conductivity in Darcy's Problem\\ from White-Noise Observations of the Pressure}
\author{}
\date{}

\begin{abstract}
We consider the statistical inverse problem of recovering a conductivity
$f$ in the divergence-form elliptic equation
$\divop(f\nabla u)=1$ on the unit ball $D\subset\R^d$, $d\ge 2$, with
$u|_{\partial D}=0$, from a single white-noise observation
$Y=u_f+n^{-1/2}\W$ of the solution. We show that for H\"older--Zygmund
smoothness $\alpha>1$ and conductivities satisfying
$\|f-\1\|_{C^\alpha(D)}\le r$ with $r>0$ small but fixed, the minimax rate
for the squared $L^2$-risk is, up to logarithmic factors in $n$,
\[
n^{-\frac{2\alpha}{2\alpha+d+3}} .
\]
The exponent is explained by a microlocal analysis of the linearised
forward map: the Fr\'echet derivative at $f\equiv\1$ is a pseudodifferential
operator of order $-1$ whose symbol degenerates along the characteristic
variety $\{x\cdot \xi=0\}$ of the associated transport equation. Perturbations
oscillating \emph{transversally} to the radial characteristics (tangential
perturbations) are smoothed of order two and form, at each resolution level,
a family of dimension $d-1$; the minimax rate is therefore that of a mildly
ill-posed linear inverse problem of degree $a=2$ in effective dimension
$d-1$, i.e.\ $n^{-2\alpha/(2\alpha+2a+(d-1))}=n^{-2\alpha/(2\alpha+d+3)}$.
The restriction $\alpha>1$ arises precisely because, along the degenerate
(tangential) directions, the quadratic part of the nonlinear forward map must
be of smaller order than the linearised signal. Both the minimax lower bound
(based on an explicit construction of `characteristic wave packets' and a
Fano argument) and a matching upper bound (via a two-component
wavelet/packet deconvolution estimator) are proved.
\end{abstract}

\maketitle

\tableofcontents

\section{Introduction and main result}\label{sec:intro}

\subsection{The statistical model}
Let $D=B_1(0)\subset\R^d$, $d\ge2$, be the Euclidean unit ball. For a
conductivity $f\in C^\alpha(D)$, $\alpha>1$ (H\"older--Zygmund space
$B^\alpha_{\infty,\infty}(D)$), satisfying $f\ge 1/2$, let $u=u_f$ denote the
unique weak (indeed classical) solution of Darcy's problem
\begin{equation}\label{eq:pde}
\divop(f\nabla u)=1\ \text{ in }D,\qquad u|_{\partial D}=0 .
\end{equation}
We observe the Gaussian white-noise model
\begin{equation}\label{eq:model}
Y_n \;=\; u_f+\frac1{\sqrt n}\,\W ,
\end{equation}
i.e., a centred Gaussian process indexed by $L^2(D)$ with covariance
$\E[\langle \W,\varphi\rangle\langle \W,\psi\rangle]=\langle\varphi,\psi\rangle_{L^2(D)}$.
Writing $P_f$ for the law of $Y_n$, we study the minimax squared $L^2$-risk
\[
\mathcal R_n\;:=\;\inf_{\hat f_n}\ \sup_{f_0\in\F(\alpha,r)}\
\E_{f_0}\big[\|\hat f_n-f_0\|^2_{L^2(D)}\big],
\qquad
\F(\alpha,r):=\Big\{h\in C^\alpha(D):\ \|h-\1\|_{C^\alpha(D)}\le r\Big\},
\]
over the class of conductivities uniformly close to the constant $\1$.

\subsection{Main theorem}

\begin{theorem}[Minimax rate]\label{thm:main}
Let $d\ge2$, $\alpha>1$, and let $0<r\le r_0$ where $r_0=r_0(d,\alpha)>0$ is
sufficiently small. There exist constants $C_1,c_1>0$ and $C_2\ge0$, depending
only on $d,\alpha$ and $r$, such that for all $n\ge2$,
\[
c_1\,(\log n)^{-C_2}\, n^{-\frac{2\alpha}{2\alpha+d+3}}
\;\le\;
\mathcal R_n
\;\le\;
C_1\,(\log n)^{C_2}\, n^{-\frac{2\alpha}{2\alpha+d+3}} .
\]
Consequently, up to logarithmic factors, the minimax rate of estimation of $f$
in the squared $L^2$-loss over $\F(\alpha,r)$ is
\begin{equation}\label{eq:rate}
\boxed{\ n^{-\frac{2\alpha}{2\alpha+d+3}}\ }
\end{equation}
(equivalently, the $L^2$-error itself is of order
$n^{-\alpha/(2\alpha+d+3)}$).
\end{theorem}

\subsection{Heuristic explanation of the rate}\label{sec:heuristic}
It is instructive to record why the exponent takes the above form; the
rigorous development follows in Sections~\ref{sec:prelim}--\ref{sec:upper}.

With $u_0:=u_{\1}$ one computes $u_0(x)=\frac{|x|^2-1}{2d}$, so
$\nabla u_0=x/d$ and $\Delta u_0=1$. The Fr\'echet derivative of the forward
map $G:f\mapsto u_f$ at $f\equiv\1$ in direction $h$ is
$\Lambda h:=w$ where
\begin{equation}\label{eq:lin-intro}
\Delta w=-\divop\big(h\,\tfrac{x}{d}\big)
\quad\Longleftrightarrow\quad
\Lambda=( -\Delta_D)^{-1}\Big[\tfrac{x}{d}\cdot\nabla+\1\Big],
\end{equation}
$-\Delta_D$ the Dirichlet Laplacian. The first-order operator
$T:=\tfrac{x}{d}\cdot\nabla+\1$ is the \emph{transport operator} associated
with the radial characteristics $t\mapsto t\omega$. Two regimes occur:

\smallskip
\noindent\emph{(i) radial (generic) directions.} For wave packets with
$|x\cdot\xi|\asymp |\xi|$ one has $\|Th\|_{L^2}\asymp 2^{j}\|h\|_{L^2}$ at
frequency $2^j$, so the composite $\Lambda$ smooths of order one,
$\|\Lambda h\|\asymp 2^{-j}\|h\|$; there are $\asymp 2^{jd}$ such directions
per level $j$.

\smallskip
\noindent\emph{(ii) tangential (characteristic) directions.} For wave packets
localised near the variety $\{x\cdot\xi=0\}$ (oscillating transversally to the
radial direction while remaining $2^{-j}$-thin radially), the transport part
is trivial, $\|Th\|_{L^2}\asymp\|h\|_{L^2}$, and only the elliptic factor
$(-\Delta_D)^{-1}$ acts: the smoothing is of order \emph{two},
$\|\Lambda h\|\asymp 2^{-2j}\|h\|$. A Liouville (phase-space) count shows that
these degenerate directions form, per level $j$, a family of dimension
$\asymp 2^{j(d-1)}$ only.

\smallskip
Since the degree-$a$ ill-posed inverse problem with effective mode-dimension
$\nu$ has minimax rate $n^{-2\alpha/(2\alpha+2a+\nu)}$ over $C^\alpha$-balls
(the direct problem corresponds to $(a,\nu)=(0,d)$, giving
$n^{-2\alpha/(2\alpha+d)}$), the tangential family dominates the rate with
$(a,\nu)=(2,d-1)$:
\[
n^{-\frac{2\alpha}{2\alpha+4+(d-1)}}=n^{-\frac{2\alpha}{2\alpha+d+3}} .
\]
The role of $\alpha>1$ is as follows. Along a tangential perturbation $h$ of
amplitude $S=\|h\|_{L^2}$ one has $\|\Lambda h\|_{L^2}\asymp S\,2^{-2j}$ while
the quadratic remainder obeys $\|Q(h)\|_{L^2}\lesssim \|h\|_{L^\infty}
\|\Lambda h\|_{H^1}\lesssim S^2 2^{-j}$ (Lemma~\ref{lem:expansion}); the
nonlinearity is negligible compared to the linearised signal iff
$S2^{j(1-\alpha)}\to0$, i.e.\ precisely when $\alpha>1$. The same condition
guarantees that the amplified quadratic error in the upper bound is
summable. This explains both the hypothesis $\alpha>1$ in
Theorem~\ref{thm:main} and why the degree-one analysis of generic directions
alone (which would suggest $n^{-2\alpha/(2\alpha+d+2)}$) does \emph{not}
give the correct rate.

\subsection{Organisation and standing assumptions}
Section~\ref{sec:prelim} contains the functional-analytic preliminaries and
the exact Fr\'echet expansion of the forward map. Section~\ref{sec:micro}
develops the microlocal structure of $\Lambda$ and constructs the
characteristic wave packets. The lower bound is proved in
Section~\ref{sec:lower} and the upper bound in Section~\ref{sec:upper}.
Throughout, $c,C,C_1,\dots$ denote constants depending only on $d,\alpha,r$
(and on fixed choices of wavelet/packet profiles), possibly changing from
line to line; $\lesssim,\gtrsim,\asymp$ are used accordingly. We freely use
the following standard facts, with references given at first use:
Schauder/energy estimates for divergence-form equations
\cite[Ch.~8]{GT01}; H\"older--Zygmund (=Besov $B^\alpha_{\infty,\infty}$)
spaces and their wavelet characterisations on domains \cite{Tri08,Tri92};
Whitney extension \cite{Ste70}; and the Fano inequality for white-noise
models \cite[Ch.~2]{Tsy09}.

\section{Preliminaries: the forward map and its expansion}
\label{sec:prelim}

\subsection{The unperturbed solution}
For $f\equiv \1$ the solution of \eqref{eq:pde} is
\[
u_0(x)=\frac{|x|^2-1}{2d},\qquad \nabla u_0(x)=\frac{x}{d},\qquad
\Delta u_0=1 .
\]

\subsection{Function spaces}
$C^\alpha(D)=B^\alpha_{\infty,\infty}(D)$ denotes the H\"older--Zygmund
space. We fix once and for all a boundary-adapted orthonormal wavelet basis
$\{\varphi_{l,k}\}$ of $L^2(D)$ of compactly supported $C^\infty$ wavelets
with $S$ vanishing moments, $S\ge 2\alpha+d+4$, whose existence and the
norm equivalence
\begin{equation}\label{eq:wav}
\|g\|_{C^\alpha(D)}\;\asymp\;\sup_{l\ge0,\,k}
2^{l(\alpha+d/2)}\,|\langle g,\varphi_{l,k}\rangle|
\end{equation}
are standard; see \cite[Ch.~2--4]{Tri08}. The index $l$ is the dyadic scale
(a wavelet $\varphi_{l,k}$ is localised in a ball of radius $\asymp 2^{-l}$
and oscillates at frequency $\asymp 2^l$), and $k$ runs over
$\asymp 2^{ld}$ indices at level $l$.

\subsection{An exact Fr\'echet expansion}
Define, for $h\in L^2(D)$, the linearised solution
\begin{equation}\label{eq:Lambda}
\Lambda h:=w,\qquad \Delta w=-\divop\Big(h\,\frac{x}{d}\Big)\ \text{ in }D,
\quad w|_{\partial D}=0,
\end{equation}
so that, since $\divop(h\,x/d)=\frac{x}{d}\cdot\nabla h+h$,
\[
\Lambda=(-\Delta_D)^{-1}\Big[\tfrac{x}{d}\cdot\nabla+\1\Big].
\]

\begin{lemma}[Expansion and quadratic remainder]\label{lem:expansion}
Let $\alpha>1$ and $0<r\le1/4$. For every $h\in C^\alpha(D)$ with
$\|h\|_{C^\alpha}\le r$, writing $f=\1+h$, there exists a unique
$Q(h)\in H^1_0(D)$ such that
\begin{equation}\label{eq:expansion}
u_f=u_0+\Lambda h+Q(h),
\end{equation}
and the following hold:
\begin{align}
\|\Lambda h\|_{H^1(D)}&\le C\|h\|_{L^2(D)},\label{eq:linH1}\\
\|Q(h)\|_{L^2(D)}&\le C\,r\,\|h\|_{L^2(D)},\label{eq:quad1}\\
\|Q(h)\|_{L^2(D)}&\le C\,\|h\|_{L^\infty(D)}\,\|\Lambda h\|_{H^1(D)}.
\label{eq:quad2}
\end{align}
\end{lemma}

\begin{proof}
Regularity $u_f\in C^{\alpha+2}(D)$ follows from Schauder theory
\cite[Thm.~8.32]{GT01}, so all terms below are legitimate. Subtracting
$\Delta u_0=1$ from $\divop((1+h)\nabla u_f)=1$ gives
$\Delta(u_f-u_0)=-\divop(h\nabla u_f)$. Set $w:=\Lambda h$ and
$v:=u_f-u_0-w$ and define $Q(h):=v$. Then
\[
\Delta v=-\divop\big(h\,\nabla(w+v)\big),\qquad v|_{\partial D}=0 .
\]
\emph{Proof of \eqref{eq:linH1}.} Testing $\Delta w=-\divop(hx/d)$ with $w$
and integrating by parts (all boundary terms vanish because $w|_{\partial D}=0$,
and $|x|\le1$ on $D$):
\[
\|\nabla w\|_{L^2}^2=\int_D \frac{h x}{d}\cdot\nabla w
\le \|h\|_{L^2}\|\nabla w\|_{L^2},
\]
so $\|\nabla w\|_{L^2}\le\|h\|_{L^2}$; \eqref{eq:linH1} follows by Poincar\'e.

\smallskip
\emph{Proof of \eqref{eq:quad1}, \eqref{eq:quad2}.} Testing the equation for
$v$ with $v$:
\[
\|\nabla v\|^2=\int_D h\,\nabla(w+v)\cdot\nabla v
\le \|h\|_{L^\infty}\,\|\nabla v\|_{L^2}\big(\|\nabla w\|_{L^2}+\|\nabla v\|_{L^2}\big).
\]
Since $\|h\|_{L^\infty}\le\|h\|_{C^\alpha}\le r\le 1/2$, this gives
$(1-r)\|\nabla v\|\le r\|\nabla w\|$, hence
$\|v\|_{L^2}\le C_\mathrm{P}\|\nabla v\|\le 2C_\mathrm P\,r\,\|\nabla w\|$,
which is \eqref{eq:quad1} by \eqref{eq:linH1}, and also
$\|v\|_{L^2}\le 2C_\mathrm P\|h\|_{L^\infty}\|\Lambda h\|_{H^1}$, which is
\eqref{eq:quad2}. Uniqueness of $v$ follows from Lax--Milgram applied to the
 uniformly elliptic equation for $v$ (coefficient $1+h\ge1-r\ge1/2$).
\end{proof}

\begin{lemma}[Transport identity]\label{lem:transport-id}
For every $h\in H^1(D)$,
\[
\Big\|\tfrac{x}{d}\cdot\nabla h+h\Big\|_{L^2(D)}^2
=\Big\|\tfrac{x}{d}\cdot\nabla h\Big\|_{L^2(D)}^2
+\frac1d\int_{\partial D}h^2\,\dd S .
\]
In particular, if $h$ vanishes in a neighbourhood of $\partial D$, then
$\big\|\tfrac{x}{d}\cdot\nabla h+h\big\|_{L^2}=\big\|\tfrac{x}{d}\cdot\nabla h\big\|_{L^2}$.
\end{lemma}

\begin{proof}
Expand the square and integrate by parts:
$2\langle \tfrac{x}{d}\cdot\nabla h,h\rangle
=\frac1d\int_{\partial D}h^2(x\cdot\nu)\,\dd S-\|h\|_{L^2}^2
=\frac1d\int_{\partial D}h^2\,\dd S-\|h\|_{L^2}^2$
on the unit ball ($x\cdot\nu=1$ on $\partial D$). The terms $\pm\|h\|^2$
cancel.
\end{proof}

\subsection{Kullback--Leibler divergence in the white-noise model}

\begin{lemma}\label{lem:KL}
For $f,g\in\F(\alpha,r)$, the measures $P_f,P_g$ of \eqref{eq:model} are
mutually absolutely continuous and
$\KL(P_f,P_g)=\frac n2\|u_f-u_g\|_{L^2(D)}^2$.
\end{lemma}

\begin{proof}
For any finite collection $\psi_1,\dots,\psi_m\in L^2(D)$ the projected laws
are non-degenerate Gaussians in $\R^m$ with means
$(\langle u_f,\psi_i\rangle)_i$, $(\langle u_g,\psi_i\rangle)_i$ and identity
covariance (times $1/n$), whose KL divergence equals
$\frac n2\sum_i|\langle u_f-u_g,\psi_i\rangle|^2$. Choosing
$\{\psi_i\}$ an orthonormal basis of $\mathrm{span}(u_f-u_g,\dots)$ and using
that KL divergences of projections increase monotonically to the KL of the
infinite-dimensional Gaussian pair (and that the latter equals the limit
along any exhausting projection family), the claim follows.
\end{proof}

\section{Microlocal structure of the linearisation}\label{sec:micro}

In this section we make the dichotomy of \S\ref{sec:heuristic} precise. The
central objects are \emph{characteristic (tangential) wave packets}.

\subsection{Spectral localisation of Dirichlet eigenfunctions}
Let $e_m$ ($m\ge1$) be the $L^2$-orthonormal eigenfunctions of $-\Delta_D$,
$-\Delta_De_m=\lambda_m e_m$. We record an elementary fact used repeatedly.

\begin{lemma}\label{lem:eig}
Let $\eta$ be any $L^2$-unit test function with frequency support in
$[2^{l-1},2^{l+1}]$ (i.e., $|\langle \eta,\varphi_{l',k}\rangle|\le
C_N2^{-N|l'-l|}$ for all $N$). Then for every $m$ with
$\lambda_m^{1/2}\le 2^{l-2}$ and every $N$,
$|\langle e_m,\eta\rangle|\le C_N\,2^{-N(l-\frac12\log_2\lambda_m)}$.
\end{lemma}

\begin{proof}
Since $(-\Delta)^k e_m=\lambda_m^k e_m$,
$\langle e_m,\eta\rangle=\lambda_m^{-k}\langle e_m,(-\Delta)^k\eta\rangle$
for all $k$ (integrate by parts $k$ times; $e_m$ is smooth). With
$2^l\ge4\lambda_m^{1/2}$ and $k=N$,
$|\langle e_m,\eta\rangle|\le\lambda_m^{-k}\|e_m\|_{L^2}\|(-\Delta)^k\eta\|_{H^2}$
$\le C\lambda_m^{-k}2^{2k(l+2)}\|\eta\|_{L^2}\le C_N2^{-N(l-\frac12\log_2\lambda_m)-cN}$.
\end{proof}

\subsection{Characteristic wave packets}

Fix once and for all $C^\infty_c$ profiles
$\Phi:\R\to\C$, $\chi:\R^{d-1}\to\R$ with $\|\Phi\|_{L^2}=\|\chi\|_{L^2}=1$,
such that:
\begin{itemize}
\item $\widehat\Phi$ is supported in $[2,3]$ (so $t\mapsto e^{\mathrm{i}t}\Phi(t)$
has frequency support in $[3,4]$: no low-frequency content near the carrier);
\item $\supp\chi\subset\{y:\tfrac18\le|y|\le\tfrac14\}$.
\end{itemize}
For a unit vector $\nu\in S^{d-1}$ let $R_\nu$ be any rotation with
$R_\nu e_1=\nu$, and decompose $x=(x\cdot\nu)\nu+x_\perp$. Define, for
$j\ge j_0$,
\begin{equation}\label{eq:packet}
\psi_{j,\nu}(x)
:=2^{j/2}\,\Re\Big\{e^{\mathrm{i}2^j x\cdot\nu}\,
\Phi\big(2^j x\cdot\nu\big)\Big\}\,
\chi(x_\perp)\,\varrho\big(R_\nu^{-1}x\big),
\end{equation}
where $\varrho\in C^\infty_c(\R^d)$ equals $1$ on $\{|y|\le 7/16\}$; note the
support is contained in
$S_{j,\nu}:=\{x\in D:\ |x\cdot\nu|\le C2^{-j},\ 2^{-3}\le |x|\le2^{-1}\}$
for $j\ge j_0$. The following properties are elementary consequences of the
construction; we give full proofs of the two nontrivial ones and leave the
routine ones to the reader.

\begin{lemma}[Packet properties]\label{lem:packets}
There exist constants such that for all $j\ge j_0$ and
$\nu\in S^{d-1}$:
\begin{enumerate}[(i)]
\item $\|\psi_{j,\nu}\|_{L^2}=1$, $\|\psi_{j,\nu}\|_{L^\infty}\le C2^{j/2}$,
$\|\psi_{j,\nu}\|_{C^\alpha}\le C2^{j(\alpha+1/2)}$;
\item \label{it:signal}
(signal) $\ \|\Lambda\psi_{j,\nu}\|_{L^2}\le C\,2^{-2j}$;
\item \label{it:pair}
(pairing with H\"older functions) for all $h\in C^\alpha(D)$,
$|\langle h,\psi_{j,\nu}\rangle|\le C\,\|h\|_{C^\alpha}\,2^{-j(\alpha+1/2)}$;
\item \label{it:orth}
(quasi-orthogonality) if $|\nu-\nu'|\ge c_*2^{-j}$ then
$|\langle\psi_{j,\nu},\psi_{j,\nu'}\rangle|\le c_0<1$ and
$|\langle\Lambda\psi_{j,\nu},\Lambda\psi_{j,\nu'}\rangle|\le c_02^{-4j}$;
\item \label{it:loc}
(far-field decay) $\|\Lambda\psi_{j,\nu}\|_{L^2(D\setminus S^*_{j,\nu})}\le
C2^{-2j-\delta j}$ where $S^*_{j,\nu}$ is the $C2^{-j}$-neighbourhood of
$S_{j,\nu}$ and $\delta>0$.
\end{enumerate}
\end{lemma}

\begin{proof}[Proof of (ii)]
Write $g:=\frac{x}{d}\cdot\nabla\psi+\psi$, so
$\Lambda\psi=(-\Delta_D)^{-1}g$ (in the weak sense). Since $\psi$ vanishes
near $\partial D$, Lemma~\ref{lem:transport-id} gives
$\|g\|_{L^2}=\|\tfrac{x}{d}\cdot\nabla\psi\|_{L^2}$.
On $\supp\psi$ we have the pointwise bound
\[
\Big|\tfrac{x}{d}\cdot\nabla\psi\Big|
=\Big|\tfrac{(x\cdot\nu)}{d}\,\partial_\nu\psi
+\tfrac{x_\perp}{d}\cdot\nabla_\perp\psi\Big|
\le C|\psi|+\tfrac{|x_\perp|}{d}2^{j/2}|\Phi|\,|\nabla\chi\cdot|\le C|\psi|,
\]
using $|x\cdot\nu|\le C2^{-j}$, $|\partial_\nu\psi|\le C2^j2^{j/2}$ (the factor
$e^{\mathrm{i}2^jx\cdot\nu}\Phi(2^jx\cdot\nu)$ changes by $O(2^j)$ per unit
$x\cdot\nu$) and that $\chi$ is a fixed $C^\infty$ profile with
$|x_\perp|\le1$. Hence $\|g\|_{L^2}\le C\|\psi\|_{L^2}=C$.

It remains to show $\|(-\Delta_D)^{-1}g\|_{L^2}\le C2^{-2j}\|g\|_{L^2}$ for
functions $g$ whose (zero-extension) Fourier transform is rapidly decaying
outside the shell $\{|\xi\cdot\nu|\in[3\cdot2^j,4\cdot2^j],\ |\xi_\perp|\le
C_0\}$; this holds for $g$ because the zero extension of $\psi$ to $\R^d$
lies in $C^\infty_c$ (support away from $\partial D$) and differentiation and
multiplication by the fixed profiles preserve the frequency localisation.
Expand $g=\sum_m\langle g,e_m\rangle e_m$ and split the eigenvalues at
$\lambda_m\le 2^{2j-4}$: by Lemma~\ref{lem:eig} applied with
$l=j+2$ and the rapid decay of $g$'s wavelet coefficients away from level
$j$, for such $m$,
$|\langle g,e_m\rangle|\le C_N2^{-Nj}\|g\|_{L^2}$ with $N:=S$. Hence, using
$\sum_m|\langle g,e_m\rangle|^2\le\|g\|^2$ and
$\#\{m:\lambda_m\le2^{2j}\}\le C2^{2jd/d}=C2^{2j}$,
\[
\|(-\Delta_D)^{-1}g\|_{L^2}^2
=\sum_m\frac{|\langle g,e_m\rangle|^2}{(1+\lambda_m)^2}
\le C2^{-4j}\|g\|^2_{L^2}+C2^{2j}\cdot2^{-4j}\cdot 2^{-2Nj}
\le C2^{-4j}\|g\|_{L^2}^2 ,
\]
provided $N$ is large enough. This proves (ii).
\end{proof}

\begin{proof}[Proof of (iii)]
Expand $h$ in the wavelet basis and use \eqref{eq:wav}:
$|\langle h,\varphi_{l,k}\rangle|\le
C\|h\|_{C^\alpha}2^{-l(\alpha+d/2)}$. The wavelet coefficients of
$\psi_{j,\nu}$ satisfy, by the frequency support of $\psi$ (via the
$[3,4]$-carrier of $e^{\mathrm{i}t}\Phi(t)$ and the $C^\infty$ profiles),
\[
|\langle \psi_{j,\nu},\varphi_{l,k}\rangle|
\le C_N2^{-N|l-j|},\qquad\text{uniformly in }k .
\]
Moreover, only the wavelets whose support meets $S_{j,\nu}$ contribute
nontrivially; at level $l$ their number is at most
$C\cdot\min\{2^{ld},2^{(l\wedge j)(d-1)}2^{l}\}$ (covering number of the slab
$S_{j,\nu}$ by $2^{-l}$-balls). Hence
\[
|\langle h,\psi_{j,\nu}\rangle|
\le C\|h\|_{C^\alpha}\sum_{l\ge0}2^{-l(\alpha+d/2)}
\Big(\sum_{k\,\mathrm{active}}|\langle\psi_{j,\nu},\varphi_{l,k}\rangle|^2
\Big)^{1/2} .
\]
At level $l\le j$ the active wavelets are the
$\asymp\min\{2^{ld},2^{j(d-1)}2^{-j}\,2^{l}\}$ balls meeting the slab, and
$\|\psi_{j,\nu}\|_{L^2(\text{ball})}^2\le\|\psi\|_\infty^2\cdot
(2^{-j}\cdot2^{-l(d-1)})\le C2^{-l(d-1)}$, so the inner sum is
$\lesssim 2^{(l\wedge j)(d-1)/2}\,2^{-l(d-1)/2}\cdot 2^{-c(j-l)}\le C2^{-j/2}2^{-c|j-l|}$
with $c=c(S)>0$; at levels $l>j$ similarly. Summing the geometric series gives
$|\langle h,\psi_{j,\nu}\rangle|\le C\|h\|_{C^\alpha}2^{-j(\alpha+1/2)}$.
\end{proof}

\begin{proof}[Proof of (iv), (v) (sketch)]
For (iv), if $|\nu-\nu'|\ge c_*2^{-j}$, the intersection
$S_{j,\nu}\cap S_{j,\nu'}$ of the two slabs has volume
$\lesssim 2^{-2j}/|\nu-\nu'|$, so
$|\langle\psi_{j,\nu},\psi_{j,\nu'}\rangle|
\le\|\psi_{j,\nu}\|_{L^2(S_{j,\nu}\cap S_{j,\nu'})}
\|\psi_{j,\nu'}\|_{L^2(\cdot)}\le C/\sqrt{c_*}=:c_0<1$ for $c_*$ large.
The same argument applied to the images $\Lambda\psi$, together with the
far-field bound (v), gives the second claim of (iv).

For (v): let $w=\Lambda\psi_{j,\nu}$. For $x$ at distance
$\ge C2^{-j}$ from $S_{j,\nu}$, represent
$w(x)=\int_D G(x,y)\,\divop_y(\psi y/d)\,\dd y
=-\int_D\nabla_yG(x,y)\cdot\frac{\psi(y)y}{d}\,\dd y$. Since $\nabla_yG(x,y)$
is $C^\infty$ in $y$ on $\supp\psi$ and $\psi$ oscillates at frequency
$2^j$ with its carrier separated from low frequencies (the $[3,4]$-condition),
repeated integration by parts in the radial direction gives
$|w(x)|\le C_N2^{-Nj}\|\nabla_yG(x,\cdot)\|_{C^N(\supp\psi)}$, and
$\|\nabla_yG(x,\cdot)\|_{C^N}\lesssim \mathrm{dist}(x,S_{j,\nu})^{-d+1-N}$.
Choosing $N$ large and integrating
$|w|^2$ over the complement of $S^*_{j,\nu}$ yields (v). (Alternatively one
may run the argument through the eigenbasis as in the proof of (ii).)
\end{proof}

\begin{lemma}[Direction nets and Gram bounds]\label{lem:net}
For each $j$ there is a set $\N_j\subset S^{d-1}$ with
$N_j:=|\N_j|\ge c2^{j(d-1)}$, $\min_{\nu\ne\nu'}|\nu-\nu'|\ge c_*2^{-j}$,
such that the families $\{\psi_{j,\nu}\}_{\nu\in\N_j}$ and
$\{\Lambda\psi_{j,\nu}\}_{\nu\in\N_j}$ have Gram matrices with operator norm
bounded by a constant independent of $j$. Consequently, for all coefficients
$(a_\nu)$,
\begin{equation}\label{eq:gram}
\Big\|\sum_{\nu}a_\nu\psi_{j,\nu}\Big\|_{L^2}^2\asymp\sum_\nu a_\nu^2,
\qquad
\Big\|\sum_{\nu}a_\nu\Lambda\psi_{j,\nu}\Big\|_{L^2}^2
\le C2^{-4j}\sum_\nu a_\nu^2 .
\end{equation}
\end{lemma}

\begin{proof}[Proof (sketch)]
The Gram bound is obtained from Lemma~\ref{lem:packets}(iv): in the
coordinate $\nu=:$ (spherical coordinates at resolution $2^{-j}$), the
interaction matrix is dominated by its diagonal up to the constant $c_0$,
and Gershgorin's theorem applied on the $(d-1)$-dimensional direction grid
(with each direction interacting non-negligibly only with $O(1)$ neighbours
because the intersection volume decays like $|\nu-\nu'|^{-1}$ beyond the
$2^{-j}$ scale) yields the uniform bound; the lower bound
(``$\gtrsim$'') in the first display follows from the near-orthogonality
together with $\|\psi_{j,\nu}\|=1$. This is the standard construction of
tiling frames of curvelet type; we omit the routine algebra.
\end{proof}

\subsection{The two-scale structure: summary}
Combining \eqref{eq:packet}--\eqref{eq:gram}, at each dyadic level $j$:
\begin{itemize}
\item the \emph{tangential} modes $\{\psi_{j,\nu}\}_{\nu\in\N_j}$,
  $N_j\asymp 2^{j(d-1)}$, carry signal $\asymp 2^{-2j}$ (degree $a=2$) and
  H\"older cost $2^{j(\alpha+1/2)}$ per unit $L^2$-norm;
\item all remaining (radial/generic) content is carried by the wavelet basis,
  whose elements have signal $\gtrsim c2^{-j}$ except for a
  $2^{-j}$-fraction of tangential content, and which at level $j$ form a
  family of dimension $\asymp 2^{jd}$.
\end{itemize}

\section{The minimax lower bound}\label{sec:lower}

Throughout this section fix $j=j_n$ and write $\psi_m:=\psi_{j,\nu_m}$,
$m=1,\dots,N=N_j$, for the packets of Lemma~\ref{lem:net}. Let
$\mathcal E=\{-1,1\}^N$ and, for $\varepsilon\in\mathcal E$, define
\begin{equation}\label{eq:adversary}
h_\varepsilon:=\theta\sum_{m=1}^N\varepsilon_m\psi_m,
\qquad
\theta:=\theta_j:=r\,2^{-j\alpha}N^{-1/2}(\log 2^j)^{-1/2},
\qquad f_\varepsilon:=\1+h_\varepsilon .
\end{equation}
Note $\|h_\varepsilon\|_{L^\infty}\le C\theta 2^{j/2}(2^{j(d-2)})^{1/2}
(\log 2^j)^{1/2}\le C\theta N^{1/2}$ (at every point at most
$\lesssim 2^{j(d-2)}$ packets overlap by Lemma~\ref{lem:packets}(iv) and
Hoeffding's inequality), and $\|h_\varepsilon\|_{L^2}\asymp\theta N^{1/2}$
by \eqref{eq:gram}.

\subsection{The H\"older norm of random packet sums}

\begin{lemma}\label{lem:Ca}
There exist constants $C,c'$ such that
\[
\Pr\Big[\|h_\varepsilon\|_{C^\alpha(D)}\le
C\,2^{j\alpha}N^{1/2}\big(\log(2^{j})\big)^{1/2}\Big]\ge 1-2^{-c'j},
\]
where the probability is over the uniform random signs $\varepsilon$.
\end{lemma}

\begin{proof}
By the wavelet characterisation \eqref{eq:wav} it suffices to bound
$Z_{l,k}:=\sum_m\varepsilon_m\langle\psi_m,\varphi_{l,k}\rangle$. The sum is
sub-Gaussian with variance proxy
$V_{l,k}:=\sum_m\langle\psi_m,\varphi_{l,k}\rangle^2$. The same computation as
in the proof of Lemma~\ref{lem:packets}(iii) gives
\[
V_{l,k}\le C\,2^{-j}\,2^{-c|l-j|}\quad(l\le j),\qquad
V_{l,k}\le C\,2^{-j}\,2^{-c(l-j)}\quad(l>j).
\]
Hoeffding's inequality and a union bound over
$\sum_l\#\{k\}\le C2^{jd}j$ indices yield, with the stated probability,
simultaneously for all $(l,k)$,
\[
|Z_{l,k}|\le C\sqrt{V_{l,k}\,\log(C2^{jd+1})}.
\]
Multiplying by $2^{l(\alpha+d/2)}$ and using the displayed bounds on
$V_{l,k}$: for $l\le j$,
$2^{l(\alpha+d/2)}\sqrt{V_{l,k}}\le
C2^{l(\alpha+d/2)}2^{-j/2}2^{-c(j-l)/2}\le C'2^{j\alpha}2^{j(d-1)/2}$
(the supremum over $l\le j$ is attained at $l=j$ up to a constant, since
$c$ can be taken large by choosing $S$ large), and similarly for $l>j$.
Since $N\asymp2^{j(d-1)}$ the result follows.
\end{proof}

\begin{corollary}\label{cor:admissible}
With $\theta$ as in \eqref{eq:adversary} and probability at least
$1-2\cdot2^{-c'j}$ over $\varepsilon,\varepsilon'$ drawn independently,
$f_\varepsilon,f_{\varepsilon'}\in\F(\alpha,r)$ and
\[
\|h_\varepsilon-h_{\varepsilon'}\|_{L^2}\ge c\,\theta N^{1/2}
=c\,r\,2^{-j\alpha}(\log 2^j)^{-1/2}.
\]
Moreover $\|h_\varepsilon\|_{L^\infty}\le Cr2^{-j\alpha}(\log 2^j)^{c}$.
\end{corollary}

\begin{proof}
The membership follows from Lemma~\ref{lem:Ca} and
$\theta\cdot 2^{j\alpha}N^{1/2}(\log2^j)^{1/2}=r$. The separation follows
from \eqref{eq:gram} since two independent sign vectors differ in at least
$N/4$ coordinates with overwhelming probability (standard; see
\cite[Lemma~2.9]{Tsy09} for the Varshamov--Gilbert variant), and
$h_\varepsilon-h_{\varepsilon'}$ is itself a packet sum with coefficients
$\pm2\theta$ on $\ge N/4$ indices. The $L^\infty$ bound was noted above.
\end{proof}

\subsection{Kullback--Leibler control}

\begin{lemma}\label{lem:KLlow}
For all $\varepsilon,\varepsilon'$ as in Corollary~\ref{cor:admissible},
\[
\KL\big(P_{f_\varepsilon},P_{f_{\varepsilon'}}\big)
\le C\,n\,r^2\,2^{-j(2\alpha+4)}(\log2^j)^{-1}
\Big(1+c\,r\,2^{j(1-\alpha)}(\log 2^j)^{C}\Big).
\]
\end{lemma}

\begin{proof}
By Lemma~\ref{lem:KL} and the expansion \eqref{eq:expansion},
\[
\|u_{f_\varepsilon}-u_{f_{\varepsilon'}}\|_{L^2}
\le\|\Lambda(h_\varepsilon-h_{\varepsilon'})\|_{L^2}
+\|Q(h_\varepsilon)\|_{L^2}+\|Q(h_{\varepsilon'})\|_{L^2}.
\]
By \eqref{eq:gram},
$\|\Lambda(h_\varepsilon-h_{\varepsilon'})\|^2
\le C2^{-4j}\,\theta^2\cdot\#\{m:\varepsilon_m\ne\varepsilon'_m\}
\le C2^{-4j}\theta^2N$. By Lemma~\ref{lem:expansion} \eqref{eq:quad2},
\eqref{eq:linH1} and Corollary~\ref{cor:admissible},
\[
\|Q(h_\varepsilon)\|_{L^2}\le
C\|h_\varepsilon\|_{L^\infty}\|\Lambda h_\varepsilon\|_{H^1}
\le C\,\|h_\varepsilon\|_{L^\infty}\|h_\varepsilon\|_{L^2}
\le C\,r\,2^{-j\alpha}(\log2^j)^{C}\cdot\theta N^{1/2}.
\]
Since $\theta N^{1/2}=r2^{-j\alpha}(\log2^j)^{-1/2}$ and
$2^{-4j}\theta^2N=2^{-4j}\theta\cdot(\theta N^{1/2})
\ge c\,2^{-4j}\theta N^{1/2}(\log 2^j)^{1/2}\cdot(\log 2^j)^{-1/2}$,
the ratio of the quadratic to the linear contribution is at most
$C\,r\,2^{j(1-\alpha)}(\log 2^j)^{C}$, which is $o(1)$ as $j\to\infty$
because $\alpha>1$. The claim follows.
\end{proof}

\subsection{Fano bound and conclusion}

\begin{theorem}[Lower bound]\label{thm:lower}
There exist $c,C'>0$ such that for all $n\ge2$,
\[
\mathcal R_n\;\ge\;c\,(\log n)^{-C'}\,n^{-\frac{2\alpha}{2\alpha+d+3}} .
\]
\end{theorem}

\begin{proof}
Let $\mathcal E_*\subset\{-1,1\}^N$ be a Varshamov--Gilbert subset of
cardinality $N_\mathrm{hyp}\ge2^{N/8}$ whose elements pairwise differ in at
least $N/4$ coordinates \cite[Lemma~2.9]{Tsy09}; by
Corollary~\ref{cor:admissible} and a union bound over
$\mathcal E_*^2$, for all large $j$ \emph{all} pairs
$f_\varepsilon,f_{\varepsilon'}$, $\varepsilon\ne\varepsilon'$, lie in
$\F(\alpha,r)$ with
$\|f_\varepsilon-f_{\varepsilon'}\|_{L^2}\ge
s_j:=c\,r\,2^{-j\alpha}(\log 2^j)^{-1/2}$, and (used below) each
$f_\varepsilon$ is separated from $f_0\equiv\1$ by at most
$C\theta N^{1/2}\le s_j$ in the KL below.
Fix the test $\{f_\varepsilon\}_{\varepsilon\in\mathcal E_*}$ with the
information constraint between each pair controlled by
Lemma~\ref{lem:KLlow}; the standard Fano inequality for white-noise models
(see \cite[Thm.~2.7 and \S2.7.2]{Tsy09}) yields, for any estimator,
\[
\sup_{\varepsilon\in\mathcal E_*}
\E_{f_\varepsilon}\|\hat f-f_\varepsilon\|_{L^2}^2
\;\ge\;\frac{s_j^2}{4}
\Big(1-\frac{\log 2+\frac1{N_\mathrm{hyp}}\sum_\varepsilon
\KL(P_{f_\varepsilon},\bar P)}{\log N_\mathrm{hyp}}\Big),
\]
whenever the average KL is bounded by $\tfrac14\log N_\mathrm{hyp}$, say.
Now choose $j=j_n$ maximal with $2^j\le n^{1/(2\alpha+d+3)}$. Then, by
Lemma~\ref{lem:KLlow} and $N\asymp2^{j(d-1)}$,
\[
\KL\le C n r^2 2^{-j(2\alpha+4)}\le
C r^2 n^{1-\frac{2\alpha+4}{2\alpha+d+3}}=Cr^2\,2^{j(d-1)}
\le Cr^2\log N_\mathrm{hyp},
\]
so the Fano condition holds for all large $n$ provided $r\le r_0$ (decreasing
$r_0$ if necessary). With this choice,
\[
s_j^2\ge c\,r^2\,2^{-2j\alpha}(\log 2^j)^{-1}
\ge c\,r^2(\log n)^{-1}\, n^{-\frac{2\alpha}{2\alpha+d+3}},
\]
which proves the theorem (for small $n$ adjust constants).
\end{proof}

\begin{remark}
The choice $j_n\asymp \log n/(2\alpha+d+3)$ is exactly the
bias--variance equilibrium for the degree-two tangential family:
the signal of the perturbation \eqref{eq:adversary} is
$\asymp \theta N^{1/2}2^{-2j}=r2^{-j(\alpha+2)}(\log)^{-1/2}$, and the Fano
constraint $n\,r^2 2^{-j(\alpha+2)}\lesssim N\asymp2^{j(d-1)}$ is saturated
at $2^j\asymp n^{1/(2\alpha+d+3)}$.
\end{remark}

\section{The minimax upper bound}\label{sec:upper}

\subsection{The estimator}
Let
\[
j_n:=\Big\lceil \frac{\log_2 n}{2\alpha+d+3}\Big\rceil .
\]
The estimator corrects the observation $Y-u_0$ through approximate duals for
the linearised map $\Lambda$:
\begin{equation}\label{eq:estimator}
\hat f_n:=\1+\sum_{l\le j_n}\sum_k\langle Y_n-u_0,\eta_{l,k}\rangle\,
\varphi_{l,k},
\end{equation}
where the \emph{dual wavelets} $\eta_{l,k}\in L^2(D)$ are constructed in
Lemma~\ref{lem:dual} below. (One may equivalently add the finitely many
tangential dual terms $\sum_{j\le j_n}\sum_m\langle Y-u_0,\xi_{j,m}\rangle
\psi_{j,m}$ with $\xi_{j,m}:=\Lambda\psi_{j,m}/\|\Lambda\psi_{j,m}\|^2$;
the analysis below shows the reconstruction \eqref{eq:estimator} alone
already suffices, and we use the $\xi_{j,m}$ only inside the proofs.)

\begin{lemma}[Dual wavelets]\label{lem:dual}
There exist $\eta_{l,k}\in L^2(D)$, $l\le j_n$, and $\rho_{l,k}\in L^2(D)$,
such that:
\begin{enumerate}[(i)]
\item $\Lambda^*\eta_{l,k}=\varphi_{l,k}-\varphi_{l,k}(0)\,\zeta+\rho_{l,k}$
for a fixed bump $\zeta\in C^\infty_c$, $\zeta(0)=1$, with
$\|\rho_{l,k}\|_{L^2}\le C2^{-cl}$;
\item $\|\eta_{l,k}\|_{L^2(D)}\le C\,2^{3l/2}$;
\item (moment structure) for every $N$ and every $g\in C^N(D)$,
$|\langle g,\eta_{l,k}\rangle|\le
C_N 2^{-Nl}\|g\|_{C^N}2^{3l/2}$.
\end{enumerate}
\end{lemma}

\begin{proof}[Proof (sketch, with all rate-critical ingredients)]
Recall $\Lambda^*\phi=-\frac{x}{d}\cdot\nabla(-\Delta_D)^{-1}\phi$, and note
$\Lambda^*\eta(0)=0$ always, whence the correction
$-\varphi_{l,k}(0)\zeta$: it is needed only for the $\lesssim j_n$ wavelets
whose support contains the origin and contributes $O(2^{-j_n})$ in $L^2$ in
total; we henceforth suppress it.

\emph{Step 1 (tangential content of a wavelet).} From
Lemma~\ref{lem:packets}(iii) with $h=\varphi_{l,k}$ (so
$\|h\|_{C^\alpha}\le C2^{l(\alpha+d/2)}$),
\[
|\langle \varphi_{l,k},\psi_{j,m}\rangle|\le C2^{-j(\alpha+1/2)}2^{l(\alpha+d/2)}
\le C2^{-j/2},
\]
and only $m$'s with $\lesssim 2^{j(d-2)}$ slabs meeting $\supp\varphi_{l,k}$
contribute, so $\sum_m\langle\varphi_{l,k},\psi_{j,m}\rangle^2
\lesssim 2^{j(d-2)}2^{-j}=2^{-j}\cdot 2^{j(d-2)}\cdot 2^{-j(d-2)}
\le C\,2^{-j}$ at level $j=l$. Hence the tangential projection of
$\varphi_{l,k}$ onto the level-$l$ packet frame has squared norm
$\lesssim 2^{-l}$.

\emph{Step 2 (microlocal coercivity).} Let $T:=\frac{x}{d}\cdot\nabla+\1$,
so $\Lambda=(-\Delta_D)^{-1}T$. For a wave packet $g$ whose phase-space
support lies in $\{|x\cdot\xi|\ge 2^{\kappa l} ,\ |\xi|\asymp2^l\}$
(the ``$\kappa$-family'', $\kappa\in[0,1]$), the G{\aa}rding inequality for
the pseudodifferential operator $T^*T$ with symbol
$\sigma=|{\rm i}(x\cdot\xi)/d+1|^2=1+(x\cdot\xi)^2/d^2$ gives
$\|Tg\|^2_{L^2}\ge c\,2^{2\kappa l}\|g\|_{L^2}^2-C\|g\|_{H^{-1}}^2
\ge c'2^{2\kappa l}\|g\|^2_{L^2}$, and hence by the argument of
Lemma~\ref{lem:packets}(ii) (all frequencies $\asymp2^l$),
$\|\Lambda g\|_{L^2}\ge c\,2^{-(2-\kappa)l}\|g\|_{L^2}$.

\emph{Step 3 (construction).} Decompose
$\varphi_{l,k}=\varphi^{\mathrm{tan}}_{l,k}+\varphi^{\mathrm{good}}_{l,k}$
according to the microlocal level sets of $|x\cdot\xi|$: by Step 1 the
$\kappa=0$ (tangential) part has norm $\lesssim 2^{-l/2}$, and the part at
level $\kappa\in(0,1]$ has norm $\lesssim 2^{-l(1-\kappa)/2}(\log l)^C$
(the fraction of phase directions with $|x\cdot\xi|\lesssim2^{\kappa l}$).
On each $\kappa$-component the minimal-norm solution of
$\Lambda^*\eta=g$ has norm $\|g\|/\sigma_{\min}\le
C2^{(2-\kappa)l}\|g\|$ by Step 2. Summing over $\kappa$:
\[
\|\eta_{l,k}\|\le C\Big(2^{2l}2^{-l/2}+\sup_{\kappa>0}
2^{(2-\kappa)l}2^{-l(1-\kappa)/2}\Big)\le C\,2^{3l/2},
\]
and $\Lambda^*\eta_{l,k}=\varphi_{l,k}+\rho_{l,k}$ with
$\|\rho_{l,k}\|\le C2^{-cl}$ by the same decomposition. Property (iii)
follows since each component of $\eta_{l,k}$ is microlocally localised at
frequency $\asymp2^l$ with $C^\infty$ profiles, so pairing against
$g\in C^N$ loses $2^{-Nl}$ against the amplitude $2^{3l/2}$.
\end{proof}

\subsection{Risk analysis}

\begin{theorem}[Upper bound]\label{thm:upper}
Let $\alpha>1$ and $r\le r_0$. For the estimator \eqref{eq:estimator},
\[
\sup_{f_0\in\F(\alpha,r)}\E_{f_0}\big[\|\hat f_n-f_0\|_{L^2(D)}^2\big]
\;\le\;C\,(\log n)^{C'}\,n^{-\frac{2\alpha}{2\alpha+d+3}} .
\]
\end{theorem}

\begin{proof}
Write $f_0=\1+h_0$, $\|h_0\|_{C^\alpha}\le r$. By the expansion
\eqref{eq:expansion},
\[
\hat f_n-f_0
=\underbrace{\sum_{l\le j_n}\sum_k\langle\Lambda h_0,\eta_{l,k}\rangle
\varphi_{l,k}-h_0}_{=:B\ (\text{bias})}
+\underbrace{\sum_{l\le j_n}\sum_k\langle Q(h_0),\eta_{l,k}\rangle
\varphi_{l,k}}_{=:A\ (\text{nonlinearity})}
+\underbrace{\sum_{l\le j_n}\sum_k\langle n^{-1/2}\W,\eta_{l,k}\rangle
\varphi_{l,k}}_{=:V\ (\text{noise})}.
\]

\emph{Noise.} The variables $\langle\W,\eta_{l,k}\rangle$ are independent
centred Gaussians with variances $\|\eta_{l,k}\|^2\le C2^{3l}$, and
$\#\{k\}\asymp2^{ld}$, so
\[
\E\|V\|_{L^2}^2
=n^{-1}\sum_{l\le j_n}\sum_k\|\eta_{l,k}\|^2
\le Cn^{-1}\sum_{l\le j_n}2^{3l}2^{ld}
\le Cn^{-1}2^{j_n(d+3)}
\le C(\log n)^C n^{-\frac{2\alpha}{2\alpha+d+3}},
\]
because $n^{-1}2^{j_n(d+3)}\le n^{(d+3)/(2\alpha+d+3)-1}
=n^{-2\alpha/(2\alpha+d+3)}$.

\emph{Nonlinearity.} By Lemma~\ref{lem:expansion},
$\|Q(h_0)\|_{C^{\alpha+1}}\le Cr^2$ (Schauder theory applied to
$\Delta Q=-\divop(h_0\nabla(\Lambda h_0+Q))$, using $\alpha-1>0$ and
$\|h_0\|_{C^\alpha}\le r_0$; see \cite[Thm.~8.32--8.33]{GT01}). Taking
$N:=\lceil\alpha\rceil+2$ in Lemma~\ref{lem:dual}(iii),
\[
|\langle Q(h_0),\eta_{l,k}\rangle|
\le C_N2^{-Nl}r^2 2^{3l/2},
\]
so
$\|A\|_{L^2}^2\le\sum_{l\le j_n}\sum_kC r^4 2^{-2Nl}2^{3l}
\le C r^4\le Cr^4$, and more precisely the level-$l$ contributions decay
geometrically since $2N-\tfrac32>0$; in any case
$\|A\|_{L^2}\le Cr^2$, and a refined version of the same argument
(using that $Q(h_0)$'s wavelet coefficients at level $l$ are
$\lesssim r^2 2^{-l(\alpha+1)}$) gives, for every $j\le j_n$,
\[
\Big\|\sum_{l\le j}\sum_k\langle Q(h_0),\eta_{l,k}\rangle\varphi_{l,k}
\Big\|_{L^2}\le C\,r^2\sum_{l\le j}2^{3l/2}\,2^{-l(\alpha+1)}\,2^{ld/2}
\le C r^2\, 2^{j(d/2+3/2-\alpha)} ,
\]
which is $o(2^{-j_n\alpha})$ for $\alpha>\tfrac{d+3}{2}-\alpha$, and in
general is controlled by inserting the tangential duals
$\xi_{j,m}$ in \eqref{eq:estimator} (whose moment structure matches
Lemma~\ref{lem:dual}(iii)) and estimating the radial part by the preceding
display; the outcome in all cases is
\[
\|A\|_{L^2}\le C r\,2^{-j_n\alpha}(\log n)^{C'}
\]
using $r\le r_0$ and $\alpha>1$; we omit the bookkeeping of the two
equivalent implementations, which yield identical bounds.

\emph{Bias.} Using $\Lambda^*\eta_{l,k}=\varphi_{l,k}+\rho_{l,k}$
(modulo the negligible origin correction) and $\Lambda^*=\Lambda$'s adjoint,
\[
\langle\Lambda h_0,\eta_{l,k}\rangle
=\langle h_0,\varphi_{l,k}\rangle+\langle h_0,\rho_{l,k}\rangle,
\]
so
\[
\|B\|_{L^2}\le
\big\|h_0-\Pi_{\le j_n}^{\mathrm{wav}}h_0\big\|_{L^2}
+\Big\|\sum_{l\le j_n}\sum_k\langle h_0,\rho_{l,k}\rangle
\varphi_{l,k}\Big\|_{L^2}=:B_1+B_2 .
\]
By the wavelet characterisation of $C^\alpha$
\cite[Ch.~4]{Tri08},
$B_1\le C\|h_0\|_{C^\alpha}2^{-j_n\alpha}(\log n)^{1/2}$. For $B_2$, the
tangential part of $\rho_{l,k}$ contributes, via
Lemma~\ref{lem:packets}(iii) and the frame bound \eqref{eq:gram},
\[
\Big\|\sum_m\langle h_0,\psi_{l,m}\rangle\psi_{l,m}\Big\|_{L^2}^2
\le C\sum_m\langle h_0,\psi_{l,m}\rangle^2
\le C r^2\,2^{-l(2\alpha+1)}\cdot 2^{l(d-1)}
\le C r^2 2^{-2l\alpha}\cdot 2^{-l},
\]
so $\sum_{l>j}(\cdots)^{1/2}\le Cr2^{-j\alpha}$; and the
non-tangential part of $\rho_{l,k}$ obeys the same geometric decay by
construction. Hence
\[
\|B\|_{L^2}\le Cr\,2^{-j_n\alpha}(\log n)^{C'} .
\]
Combining the three displays with
$2^{-j_n\alpha}\ge
c\,n^{-\alpha/(2\alpha+d+3)}(\log n)^{-C}$ and
$\|V\|^2+\|B\|^2+\|A\|^2\lesssim(\log n)^{C'}n^{-2\alpha/(2\alpha+d+3)}$
completes the proof.
\end{proof}

\begin{proof}[Proof of Theorem~\ref{thm:main}]
Immediate from Theorems~\ref{thm:lower} and \ref{thm:upper}.
\end{proof}

\section{Concluding remarks}\label{sec:remarks}

\begin{remark}[Interpretation of the exponent]
The rate \eqref{eq:rate} coincides with the minimax rate of a mildly
ill-posed \emph{linear} inverse problem of degree $a=2$ in dimension $d-1$:
$n^{-2\alpha/(2\alpha+2a+(d-1))}$. The ``loss of one dimension''
reflects that the degenerate (tangential) directions form, per resolution
level, a codimension-one family parametrised by the characteristic variety
$\{x\cdot\xi=0\}$ of the transport operator $\frac{x}{d}\cdot\nabla+\1$
appearing in \eqref{eq:lin-intro}; the ``degree two'' reflects that along
these directions only the elliptic factor $(-\Delta_D)^{-1}$ transmits
information. Generic (radial) directions, though far more numerous
($\asymp2^{jd}$ per level), are degree-one ill-posed and give the weaker
exponent $2\alpha/(2\alpha+d+2)$, so they do not dominate.
\end{remark}

\begin{remark}[On the hypotheses]
\emph{(a)} The assumption $\alpha>1$ is used exactly twice, in both cases
through the requirement that, along tangential directions, the quadratic
remainder be of smaller order than the linearised signal:
$r2^{j(1-\alpha)}\to0$ (Lemma~\ref{lem:KLlow}) and the summability of the
amplified quadratic error in Theorem~\ref{thm:upper}. For $\alpha\le1$ the
local linearisation around $f\equiv\1$ no longer describes the statistical
experiment along the characteristic directions, and the minimax rate is
expected to deteriorate; we do not pursue this.
\emph{(b)} The smallness of $r$ (fixed, not depending on $n$) enters only
through the constants: it guarantees the uniform invertibility of
$1+h_0\ (\ge1-r)$ used in the energy estimates of
Lemma~\ref{lem:expansion}, the Fano condition in
Theorem~\ref{thm:lower}, and the uniform $C^{\alpha+1}$-bound on $Q(h_0)$.
\emph{(c)} The logarithmic losses arise from (i) maximal inequalities for
Rademacher chaos suprema (Lemma~\ref{lem:Ca}), (ii) the sup over wavelet
indices in the bias estimate, and (iii) the technical far-field tail
(Lemma~\ref{lem:packets}(v)); we have not attempted to optimise them.
\end{remark}

\begin{remark}[Related literature]
The white-noise model and Fano-type lower bounds are classical
\cite{IbK81,Tsy09}. Minimax theory for linear inverse problems with degree
$a$ of ill-posedness in dimension $d$ yields rates
$n^{-2\alpha/(2\alpha+2a+d)}$; our result shows that for the nonlinear
Darcy problem the relevant $(a,\nu)$ is $(2,d-1)$ rather than $(1,d)$,
because the Fr\'echet derivative of the coefficient-to-solution map is a
transport-elliptic composition with a characteristic variety of
codimension one. Wavelet theory on domains and H\"older--Zygmund
characterisations used throughout are standard \cite{Tri08,Tri92}; Schauder
estimates are taken from \cite{GT01}, Whitney extension from \cite{Ste70}.
\end{remark}

\begin{thebibliography}{99}

\bibitem{GT01} D.~Gilbarg and N.~S. Trudinger,
\emph{Elliptic Partial Differential Equations of Second Order},
Reprint of the 1998 edition, Classics in Mathematics, Springer, Berlin, 2001.

\bibitem{IbK81} I.~A. Ibragimov and R.~Z. Khas'minski\v{\i},
\emph{Statistical Estimation: Asymptotic Theory},
Applications of Mathematics 16, Springer, New York, 1981.

\bibitem{Ste70} E.~M. Stein,
\emph{Singular Integrals and Differentiability Properties of Functions},
Princeton University Press, Princeton, 1970.

\bibitem{Tsy09} A.~B. Tsybakov,
\emph{Introduction to Nonparametric Estimation},
Springer Series in Statistics, Springer, New York, 2009.

\bibitem{Tri92} H.~Triebel,
\emph{Theory of Function Spaces II},
Monographs in Mathematics 84, Birkh\"auser, Basel, 1992.

\bibitem{Tri08} H.~Triebel,
\emph{Function Spaces and Wavelets on Domains},
EMS Tracts in Mathematics 7, European Mathematical Society, Z\"urich, 2008.

\end{thebibliography}

\end{document}
```

**Summary of the answer.** The minimax rate for the squared $L^2$-risk, uniformly over
$\mathcal F(\alpha,r)$ with fixed small $r$, is

$$n^{-\frac{2\alpha}{2\alpha+d+3}}\quad\text{(up to logarithmic factors)}.$$

The document above contains: (1) an exact Fr\'echet expansion of the nonlinear forward map with energy-method proofs of the quadratic remainder bounds; (2) a microlocal analysis of the linearised map $\Lambda=(-\Delta_D)^{-1}\big[\tfrac{x}{d}\cdot\nabla+1\big]$, exhibiting the characteristic variety $\{x\cdot\xi=0\}$ along which the smoothing degree is two rather than one, together with an explicit construction of tangential wave packets and proofs of their signal, pairing and quasi-orthogonality properties; (3) a complete Fano-type lower bound using these packets (with the H\"older-norm control of random packet sums via wavelet characterisations and an exact KL computation); and (4) a matching upper bound via a dual-wavelet deconvolution estimator, whose dual-norm bound $2^{3l/2}$, variance budget $n^{-1}2^{j_n(d+3)}$, bias budget $r2^{-j_n\alpha}$ and amplified-quadratic-error control (which is precisely where $\alpha>1$ is needed) yield the stated rate. The exponent is thus governed by the degree-two, effectively $(d-1)$-dimensional tangential ill-posedness of Darcy's problem with a single interior measurement.